\documentclass[10pt,a4paper]{article}

% DRAFT for MS3 (due Thu 03.12.2026, 23:59). TODO(author) marks what is still open; compile to
% docs/ms3_summary.pdf (max. 2 pages) once the numbers from atc.darker.dev and the MS2 reviews are in.

\usepackage[a4paper,margin=1.7cm,top=1.4cm,bottom=1.4cm]{geometry}
\usepackage[T1]{fontenc}
\usepackage[utf8]{inputenc}
\usepackage{lmodern}
\usepackage{microtype}
\usepackage{enumitem}
\usepackage{titlesec}
\usepackage[hidelinks]{hyperref}

\setlist{nosep,leftmargin=1.1em,topsep=1pt}
\titleformat{\section}{\large\bfseries}{\thesection}{0.6em}{}
\titlespacing*{\section}{0pt}{6pt}{2pt}
\setlength{\parindent}{0pt}
\setlength{\parskip}{2.5pt}
\pagestyle{empty}
\hypersetup{pdftitle={MS3 Summary - WinterLoad},pdfauthor={Björn Pfoster}}

% label, title, status, commit
\newcommand{\point}[4]{\textbf{#1 \textperiodcentered{} #2 \textperiodcentered{} #3}\if\relax\detokenize{#4}\relax\else{} (\texttt{#4})\fi.}

\begin{document}
\small

{\raggedright\Large\bfseries MS3 Summary \textperiodcentered{} WinterLoad: Day-Ahead Forecasting of Swiss Electricity Demand\par}
\vspace{2pt}
Björn Pfoster \textperiodcentered{} MLOps HS26 \textperiodcentered{} \url{https://github.com/D4rkMindz/hslu-winterload}

\section{What I achieved in this milestone}
\begin{itemize}
  \item \textbf{Every training run goes through Kestra and lands in one MLflow.} Kestra on
        \url{https://atc.darker.dev} pulls the flows from \texttt{main} and the image that CI published, and logs to
        \url{https://ml.darker.dev}. \texttt{weekly\_training} runs champion/challenger every Monday,
        \texttt{tuning} runs monthly and by hand, \texttt{final\_test} only by hand, and \texttt{daily\_monitoring}
        starts a retrain when a trigger fires. Each stage has its own experiment (\texttt{winterload/training},
        \texttt{tuning}, \texttt{monitoring}, \texttt{studies}), and every run carries the commit of its image and the
        frames it trained on as MLflow datasets. The same flows run locally with docker compose against a local MLflow.
  \item \textbf{Tuning with a gate before the test year.} Optuna searches seven LightGBM settings. Each trial scores
        six chronological 8-week folds inside TRAIN, as a nested MLflow run. The run first scores the untuned model on
        VALIDATION (5.68\%), which sets the target of 5.12\% (B1/A4 in MS2), and then scores only the best trial on
        VALIDATION, once. TODO(author): best fold MAPE, validation MAPE and whether the gate passed, from the first
        tuning run on atc (run ID).
  \item \textbf{Settings as code, not as a lookup.} A tuning run leaves a \texttt{tuned\_settings.py} artifact. I
        paste it into \path{src/winterload/forecast/model/tuned.py} and commit it, together with the run ID. Weekly
        retraining then reads the committed settings and logs that run ID, so every registered model leads back to its
        tuning run and commit, and a rollback is a \texttt{git revert}.
  \item \textbf{One read of the test year.} \texttt{final\_test} first checks in MLflow that the committed tuning run
        exists, passed the gate and holds exactly the committed settings. If MLflow is down or a check fails, the
        execution fails with one error line and the year stays closed. TODO(author): test-year MAPE against the 6.5\%
        target and the baseline's 7.6\%.
  \item \textbf{Drift triggers a retrain again.} In MS2 I dropped drift as a trigger. I put it back:
        \texttt{retrain\_monitor} now retrains when the 7-day MAPE exceeds $0.9\times$ the baseline's \emph{or} at least
        half of the monitored columns drifted (Evidently, default tests). The HTML report goes to MLflow. Before the
        final test the reference is the data the champion trained on, because the same weeks a year earlier lie in the
        closed test year. Champion/challenger still only promotes a better model, and it skips when a comparison already
        ran in the last 24 hours, so the weekly flow and the trigger can't promote twice on one day.
        TODO(author): one sentence on how often drift fired in the first weeks of the heating season.
  \item \textbf{CI.} Ruff and 160 unit tests run on every pull request and push, and the image for atc is only
        published after they passed. The tests include an end-to-end tuning run against a throwaway MLflow.
  \item \textbf{What it cost.} LightGBM starts one thread per core. On my 16-core machine, 2000 trees with 255 leaves
        took 403\,s with 16 threads and 4.4\,s with 4, so the model now uses 4 threads. The local Kestra moved from 1.1 to
        2.0 to match atc.
  \item \textbf{Open question on deployment.} atc, MLflow and the notebooks run on my own server (srv67); only the
        features live in Google Cloud (BigQuery). TODO(author): confirm with the lecturers before MS4 that this counts as
        the required cloud deployment.
  \item \textbf{Look at first:} \path{docker/kestra/files/flows/}, \path{src/winterload/pipeline/training/} (tuning,
        champion/challenger, final test), \path{src/winterload/forecast/lifecycle/}, the experiments on ml.darker.dev,
        and \path{docs/media/system_design.svg}.
\end{itemize}

\section{Response to MS2 reviews}
TODO(author): one entry per point from both reviews and the editorial summary, with status (Changed / Partly /
Rejected / Not yet) and the commit or file. Same point from both reviewers: answer once, name both labels.

\point{A1}{TODO title}{TODO status}{}
TODO.

\point{B1}{TODO title}{TODO status}{}
TODO.

\point{E1}{TODO title}{TODO status}{}
TODO.

\end{document}
